% =====================================================================
%  Manual de comisionamiento — Cuartos fríos (módulo ColdRoomPan)
%  Diseñado para lectura fácil en campo. Compilar: pdflatex main.tex (x2)
% =====================================================================
\documentclass[11pt,a4paper]{article}

\usepackage[utf8]{inputenc}
\usepackage[T1]{fontenc}
\usepackage{amssymb}
\usepackage[strings]{underscore}      % permite _ en texto (ColdRoom_1) sin escapar
\usepackage[margin=2.3cm]{geometry}
\usepackage[scaled=0.95]{helvet}
\renewcommand{\familydefault}{\sfdefault}
\usepackage{microtype}
\usepackage[table]{xcolor}
\usepackage{graphicx}
\usepackage{booktabs}
\usepackage{array}
\usepackage{titlesec}
\usepackage[most]{tcolorbox}
\usepackage{enumitem}
\usepackage{fancyhdr}
\usepackage[hidelinks]{hyperref}

% ---------- paleta ----------
\definecolor{frio}{HTML}{1274B3}
\definecolor{frioSoft}{HTML}{E4F1FA}
\definecolor{calor}{HTML}{C06A1E}
\definecolor{calorSoft}{HTML}{FAEEDD}
\definecolor{verde}{HTML}{2C8F50}
\definecolor{verdeSoft}{HTML}{E6F4EB}
\definecolor{tinta}{HTML}{15262F}
\definecolor{gris}{HTML}{54646E}
\definecolor{linea}{HTML}{C9D6DE}
\definecolor{zebra}{HTML}{EFF5F9}
\definecolor{cab}{HTML}{1274B3}

\arrayrulecolor{linea}
\setlength{\arrayrulewidth}{0.5pt}
\renewcommand{\arraystretch}{1.28}

% ---------- macros ----------
\newcommand{\ra}{\textcolor{frio}{$\rightarrow$}}
\newcommand{\md}{\,\textcolor{gris}{\textbullet}\,}
\newcommand{\gC}{\textdegree C}
\newcommand{\slot}[1]{\textbf{#1}}
\newcommand{\code}[1]{\texttt{#1}}

% ---------- títulos ----------
\titleformat{\section}{\LARGE\bfseries\color{frio}}{\thesection.}{0.5em}{}
  [\vspace{3pt}{\color{linea}\titlerule[1.2pt]}]
\titlespacing*{\section}{0pt}{20pt}{10pt}
\titleformat{\subsection}{\large\bfseries\color{tinta}}{\thesubsection}{0.5em}{}
\titlespacing*{\subsection}{0pt}{14pt}{5pt}

% ---------- cajas ----------
\newtcolorbox{importante}{enhanced,colback=frioSoft,colframe=frio,boxrule=0.6pt,
  arc=3pt,left=10pt,right=10pt,top=7pt,bottom=7pt,fonttitle=\bfseries\color{frio},
  coltitle=frio,title={\hspace{-2pt}IMPORTANTE},attach title to upper=\quad}
\newtcolorbox{ojo}{enhanced,colback=calorSoft,colframe=calor,boxrule=0.6pt,
  arc=3pt,left=10pt,right=10pt,top=7pt,bottom=7pt,fonttitle=\bfseries\color{calor},
  coltitle=calor,title={\hspace{-2pt}OJO},attach title to upper=\quad}
\newtcolorbox{nota}{enhanced,colback=gray!7,colframe=gray!45,boxrule=0.5pt,
  arc=3pt,left=10pt,right=10pt,top=6pt,bottom=6pt,fonttitle=\bfseries\color{gris},
  coltitle=gris,title={\hspace{-2pt}NOTA},attach title to upper=\quad}
% caja de comportamiento
\newtcolorbox{comporta}{enhanced,colback=verdeSoft,colframe=verde,boxrule=0.6pt,
  arc=3pt,left=10pt,right=10pt,top=7pt,bottom=7pt}
% caja de árbol (monospace)
\newtcolorbox{arbol}{enhanced,colback=gray!6,colframe=linea,boxrule=0.5pt,arc=2pt,
  left=10pt,right=10pt,top=8pt,bottom=8pt,fontupper=\ttfamily\small}
% resumen de cuarto (banda)
\newtcolorbox{ficha}[1]{enhanced,colback=frio,colframe=frio,arc=3pt,
  left=12pt,right=12pt,top=8pt,bottom=8pt,coltext=white,fontupper=\bfseries,#1}

% ---------- encabezado/pie ----------
\setlength{\headheight}{14pt}
\pagestyle{fancy}\fancyhf{}
\renewcommand{\headrulewidth}{0.4pt}
\renewcommand{\headrule}{\hbox to\headwidth{\color{linea}\leaders\hrule height \headrulewidth\hfill}}
\fancyhead[L]{\small\color{gris}Comisionamiento — Cuartos fríos}
\fancyhead[R]{\small\color{gris}\code{ColdRoomPan}}
\fancyfoot[C]{\small\color{gris}\thepage}

% ---------- estilo de tablas ----------
\newcommand{\cabfila}{\color{white}\bfseries}
\newcolumntype{L}[1]{>{\raggedright\arraybackslash}p{#1}}

\hypersetup{colorlinks=true,linkcolor=frio,urlcolor=frio}

\begin{document}

% ================= PORTADA =================
\begin{titlepage}
\centering
\vspace*{2.2cm}
{\color{frio}\rule{\linewidth}{2pt}}\\[10pt]
{\Huge\bfseries\color{tinta} Manual de comisionamiento}\\[10pt]
{\LARGE\bfseries\color{frio} Cuartos fríos}\\[6pt]
{\large\color{gris} Módulo Niagara N4 \code{ColdRoomPan}}\\[8pt]
{\color{frio}\rule{\linewidth}{2pt}}\\[26pt]

\begin{minipage}{0.82\linewidth}\color{gris}\itshape\centering
Guía práctica para el técnico de controles que arma las estaciones reales.
Cuarto por cuarto: qué arrastrar, qué va dentro de cada bloque, la numeración,
la configuración y todo el cableado de entradas y salidas.
\end{minipage}

\vfill
\begin{minipage}{0.9\linewidth}\small\color{gris}
\textbf{Antes de empezar, esto ya debe estar hecho:} el módulo firmado
\code{coldRoomPan-rt.jar} instalado en \code{!modules} de la estación, y los puntos
físicos (proxy points) ya creados bajo su dispositivo en el árbol de drivers.\\[6pt]
\textbf{Documentos relacionados:} diseño interno (\code{cold-room-module-design.md}),
cómo se construyó el módulo (\code{how-to-create-coldroom-module.md}),
flujo de build/firma/deploy (\code{module-dev-workflow.md}).
\end{minipage}
\vspace{1.4cm}
\end{titlepage}

\tableofcontents
\clearpage

% ================= 1. INTRODUCCIÓN =================
\section{Introducción y convención de nombres}

\subsection{Los tres bloques del módulo}
El módulo aporta tres tipos de componente. En Workbench los verás con su nombre en español:

\begin{center}\small
\begin{tabular}{L{3.2cm} L{3.2cm} L{7.2cm}}
\toprule
\rowcolor{cab}\cabfila \textbf{Tipo} & \cabfila \textbf{En Workbench} & \cabfila \textbf{Qué hace} \\
\midrule
\rowcolor{zebra}\code{ColdRoom} & \textbf{Cuarto frio} & Contenedor por cuarto; maneja la lógica de enfriamiento. \\
\code{EvaporatorUnit} & \textbf{Evaporadora} & Una por evaporadora física; secuencia válvula\,\ra\,evaporadora. \\
\rowcolor{zebra}\code{DefrostController} & \textbf{Control de deshielo} & Solo Cuarto 3; coordina el deshielo entre unidades. \\
\bottomrule
\end{tabular}
\end{center}

\subsection{Numeración}
\begin{itemize}[leftmargin=1.4em,itemsep=2pt]
  \item \textbf{Un \code{ColdRoom} por cuarto}: \code{ColdRoom_1}, \code{ColdRoom_2}, \code{ColdRoom_3}, \code{ColdRoom_4}.
  \item \textbf{Dentro} de cada cuarto, las evaporadoras se numeran \code{EvaporatorUnit_1}, \code{EvaporatorUnit_2}, \dots\ \emph{por cuarto} (cada cuarto arranca en \code{_1}).
  \item El \textbf{Cuarto 3} además lleva un \code{DefrostController_1} dentro del \code{ColdRoom_3}.
\end{itemize}

\subsection{El modelo de contención (leer antes de cablear)}
\begin{importante}
Las \textbf{evaporadoras} y el \textbf{control de deshielo} van \textbf{DENTRO} del \code{ColdRoom},
como hijos del bloque — no como bloques sueltos en \code{/Config}.\\[4pt]
El \code{ColdRoom} maneja a sus hijos \textbf{internamente}: por eso \textbf{NO se cablea
bloque contra bloque}. No jales un link del \code{ColdRoom} a la \code{EvaporatorUnit}; esa
relación ya existe por estar adentro.
\end{importante}

\noindent Tu único trabajo de cableado (BLink) es enlazar los \textbf{puntos físicos} con los slots:
\begin{itemize}[leftmargin=1.4em,itemsep=2pt]
  \item \textbf{Entradas} (sensores) \ra\ slots de entrada del bloque (enlace de solo lectura).
  \item \textbf{Salidas} del bloque \ra\ \textbf{writables} físicos, escribiendo en el nivel de
        prioridad \code{in8} (nivel del programa; los niveles superiores quedan para overrides manuales).
\end{itemize}

\subsection{Nombres para los puntos físicos}
Usa un patrón consistente para ubicar cualquier punto de un vistazo (\code{C}\,=\,cuarto, \code{E}\,=\,evaporadora):

\begin{center}\small
\begin{tabular}{L{6.5cm} L{6.5cm}}
\toprule
\rowcolor{cab}\cabfila \textbf{Punto} & \cabfila \textbf{Ejemplo de nombre} \\
\midrule
\rowcolor{zebra}Sensor de zona & \code{Zona1_C1}, \code{Zona2_C1} \\
Sensor de la evaporadora (serpentín) & \code{Coil_C1E1} \\
\rowcolor{zebra}Válvula (writable) & \code{Valv_C1E1} \\
Evaporadora / ventilador (writable) & \code{Evap_C1E1} \\
\rowcolor{zebra}Resistencia de deshielo (writable, Cuarto 3) & \code{Resist_C3E1} \\
Sensor de temp de resistencia (Cuarto 3, opcional) & \code{TempResist_C3E1} \\
\bottomrule
\end{tabular}
\end{center}

\subsection{Alarmas visuales}
\begin{nota}
Las alarmas (temp de evaporadora alta/baja, temp del cuarto alta, ventilador detenido) son
\textbf{solo visuales}: no intervienen en el control. Se arman como \textbf{extensiones de punto}
(\code{BAlarmSourceExt}) sobre los sensores físicos, \emph{no} como slots de estos bloques.
Los slots \slot{Limite alarma alta del cuarto}, \slot{Limite alarma alta} y \slot{Limite alarma baja}
\emph{solo guardan el umbral}. El armado completo de alarmas es un documento aparte.
\end{nota}

% ================= 2. PASOS GENERALES =================
\section{Pasos generales (todos los cuartos)}
Sigue esta secuencia en cada cuarto; el detalle específico está en las secciones 3 a 6.

\begin{enumerate}[leftmargin=1.6em,itemsep=3pt]
  \item \textbf{Crear el cuarto:} arrastra un \code{ColdRoom} (Cuarto frio) desde la paleta a
        \code{/Config} (ej. \code{/Config/CuartosFrios/}). Renómbralo \code{ColdRoom_N}.
  \item \textbf{Agregar las evaporadoras DENTRO:} arrastra una \code{EvaporatorUnit} \emph{sobre}
        el \code{ColdRoom_N} (queda como hijo). Renómbrala \code{EvaporatorUnit_1}; repite \code{_2}, \code{_3}.
  \item \textbf{(Solo Cuarto 3) Agregar el deshielo:} arrastra un \code{DefrostController} dentro del
        \code{ColdRoom_3}. Renómbralo \code{DefrostController_1}.
  \item \textbf{Fijar el \slot{Modo de etapas}:} \textbf{Por etapas} (Cuarto 1) o \textbf{Simple} (Cuartos 2, 3, 4).
  \item \textbf{Configurar} \slot{Consigna}, \slot{Diferencial arriba}, \slot{Diferencial abajo}.
  \item \textbf{Configurar cada evaporadora:} \slot{Retardo de arranque} y \slot{Tiene deshielo}
        (Sí solo en Cuarto 3). Cargar límites de alarma si aplica.
  \item \textbf{(Solo Cuarto 3) Configurar el deshielo:} \slot{Modo}, \slot{Intervalo}/\slot{Entrada de horario},
        \slot{Duracion}, \slot{Retardo de escalonamiento}, y opcional \slot{Terminar por temp resistencia} + \slot{Umbral temp resistencia}.
  \item \textbf{Enlazar los puntos físicos} a los slots (tablas de ENTRADAS y SALIDAS de cada cuarto).
        Las salidas escriben en el nivel \code{in8}.
\end{enumerate}

\begin{nota}
\textbf{Cableado masivo:} en comisionamiento conviene usar el \code{BBatchLinkEditor} con \code{checkLink}
(dry-run) para enlazar en lote y mantener los links estables por handle/tag.
\end{nota}

% ================= 3. CUARTO 1 =================
\section{Cuarto 1 — 3 evaporadoras, 2 zonas, sin deshielo}
\begin{ficha}{}3 evaporadoras \md\ 2 zonas \md\ consigna compartida \md\ Modo de etapas = Por etapas \md\ sin deshielo\end{ficha}

\subsection{Estructura de bloques}
\begin{arbol}
/Config/CuartosFrios/\\
`-- ColdRoom_1           (Cuarto frio)  -- Modo de etapas = Por etapas\\
\phantom{`-- }|-- EvaporatorUnit_1  (Evaporadora)  -- sigue Zona 1\\
\phantom{`-- }|-- EvaporatorUnit_2  (Evaporadora)  -- sigue Zona 1 O Zona 2\\
\phantom{`-- }`-- EvaporatorUnit_3  (Evaporadora)  -- sigue Zona 2
\end{arbol}

\subsection{Configuración}
\begin{center}\small
\rowcolors{2}{white}{zebra}
\begin{tabular}{L{3.6cm} L{4.6cm} L{5.4cm}}
\rowcolor{cab}\cabfila \textbf{Bloque} & \cabfila \textbf{Slot} & \cabfila \textbf{Valor} \\
\code{ColdRoom_1} & Modo de etapas & \textbf{Por etapas} \\
\code{ColdRoom_1} & Consigna & consigna del cuarto (ej. \(-18\)\,\gC) \\
\code{ColdRoom_1} & Diferencial arriba & según proyecto (ej. 1.0) \\
\code{ColdRoom_1} & Diferencial abajo & según proyecto (ej. 1.0) \\
\code{ColdRoom_1} & Limite alarma alta del cuarto & umbral (solo dato para la alarma) \\
\code{EvaporatorUnit_1..3} & Retardo de arranque & \textbf{2 s} \\
\code{EvaporatorUnit_1..3} & Tiene deshielo & \textbf{No} \\
\code{EvaporatorUnit_1..3} & Limite alarma alta / baja & umbral de temp de serpentín (solo dato) \\
\end{tabular}
\end{center}
\begin{nota}\slot{Salida resistencia} y \slot{Temp resistencia} \textbf{no se usan} en este cuarto (no hay deshielo).\end{nota}

\subsection{Entradas y salidas}
\begin{minipage}[t]{0.49\linewidth}\small
\centering\textbf{ENTRADAS} (sensor \ra\ slot)\\[4pt]
\rowcolors{2}{white}{zebra}
\begin{tabular}{L{2.5cm} L{3.9cm}}
\rowcolor{cab}\cabfila \textbf{Punto} & \cabfila \textbf{Slot destino} \\
\code{Zona1_C1} & \code{ColdRoom_1} \md \slot{Zona 1} \\
\code{Zona2_C1} & \code{ColdRoom_1} \md \slot{Zona 2} \\
\code{Coil_C1E1} & \code{EvapUnit_1} \md \slot{Temp evap.} \\
\code{Coil_C1E2} & \code{EvapUnit_2} \md \slot{Temp evap.} \\
\code{Coil_C1E3} & \code{EvapUnit_3} \md \slot{Temp evap.} \\
\end{tabular}
\end{minipage}\hfill
\begin{minipage}[t]{0.49\linewidth}\small
\centering\textbf{SALIDAS} (slot \ra\ writable, \code{in8})\\[4pt]
\rowcolors{2}{white}{zebra}
\begin{tabular}{L{3.9cm} L{2.4cm}}
\rowcolor{cab}\cabfila \textbf{Slot origen} & \cabfila \textbf{Writable} \\
\code{EvapUnit_1} \md \slot{Salida valvula} & \code{Valv_C1E1} \\
\code{EvapUnit_1} \md \slot{Salida evap.} & \code{Evap_C1E1} \\
\code{EvapUnit_2} \md \slot{Salida valvula} & \code{Valv_C1E2} \\
\code{EvapUnit_2} \md \slot{Salida evap.} & \code{Evap_C1E2} \\
\code{EvapUnit_3} \md \slot{Salida valvula} & \code{Valv_C1E3} \\
\code{EvapUnit_3} \md \slot{Salida evap.} & \code{Evap_C1E3} \\
\end{tabular}
\end{minipage}

\subsection{Cómo se comporta}
\begin{comporta}
Escalonamiento con consigna compartida:
\begin{itemize}[leftmargin=1.3em,itemsep=1pt,topsep=3pt]
  \item \textbf{Evaporadora 1} sigue a \textbf{Zona 1}.
  \item \textbf{Evaporadora 3} sigue a \textbf{Zona 2}.
  \item \textbf{Evaporadora 2} arranca con \textbf{Zona 1 O Zona 2} (cualquiera que pida frío).
\end{itemize}
En cada evaporadora, al pedir frío se \textbf{abre la válvula primero} y, tras el \slot{Retardo de arranque}
(2 s), \textbf{arranca la evaporadora}. Al parar: primero la evaporadora, luego la válvula.
\end{comporta}

% ================= 4. CUARTO 2 =================
\section{Cuarto 2 — 1 evaporadora, 1 zona, sin deshielo}
\begin{ficha}{}1 evaporadora \md\ 1 zona \md\ Modo de etapas = Simple \md\ sin deshielo\end{ficha}

\subsection{Estructura de bloques}
\begin{arbol}
/Config/CuartosFrios/\\
`-- ColdRoom_2           (Cuarto frio)  -- Modo de etapas = Simple\\
\phantom{`-- }`-- EvaporatorUnit_1  (Evaporadora)
\end{arbol}

\subsection{Configuración}
\begin{center}\small
\rowcolors{2}{white}{zebra}
\begin{tabular}{L{3.6cm} L{4.6cm} L{5.4cm}}
\rowcolor{cab}\cabfila \textbf{Bloque} & \cabfila \textbf{Slot} & \cabfila \textbf{Valor} \\
\code{ColdRoom_2} & Modo de etapas & \textbf{Simple} \\
\code{ColdRoom_2} & Consigna & consigna del cuarto \\
\code{ColdRoom_2} & Diferencial arriba / abajo & según proyecto \\
\code{ColdRoom_2} & Limite alarma alta del cuarto & umbral (solo dato) \\
\code{EvaporatorUnit_1} & Retardo de arranque & según proyecto (ej. 2 s) \\
\code{EvaporatorUnit_1} & Tiene deshielo & \textbf{No} \\
\end{tabular}
\end{center}
\begin{nota}\slot{Zona 2} no se usa (déjala sin enlazar). \slot{Salida resistencia} y \slot{Temp resistencia} tampoco.\end{nota}

\subsection{Entradas y salidas}
\begin{minipage}[t]{0.49\linewidth}\small
\centering\textbf{ENTRADAS}\\[4pt]
\rowcolors{2}{white}{zebra}
\begin{tabular}{L{2.5cm} L{3.9cm}}
\rowcolor{cab}\cabfila \textbf{Punto} & \cabfila \textbf{Slot destino} \\
\code{Zona1_C2} & \code{ColdRoom_2} \md \slot{Zona 1} \\
\code{Coil_C2E1} & \code{EvapUnit_1} \md \slot{Temp evap.} \\
\end{tabular}
\end{minipage}\hfill
\begin{minipage}[t]{0.49\linewidth}\small
\centering\textbf{SALIDAS} (\code{in8})\\[4pt]
\rowcolors{2}{white}{zebra}
\begin{tabular}{L{3.9cm} L{2.4cm}}
\rowcolor{cab}\cabfila \textbf{Slot origen} & \cabfila \textbf{Writable} \\
\code{EvapUnit_1} \md \slot{Salida valvula} & \code{Valv_C2E1} \\
\code{EvapUnit_1} \md \slot{Salida evap.} & \code{Evap_C2E1} \\
\end{tabular}
\end{minipage}

\subsection{Cómo se comporta}
\begin{comporta}
Modo simple: cuando \textbf{Zona 1} pide frío, arranca la evaporadora (abre válvula, luego evaporadora
tras el retardo). Sin escalonamiento ni segunda zona.
\end{comporta}

% ================= 5. CUARTO 3 =================
\section{Cuarto 3 — 2 evaporadoras, 1 zona, CON deshielo}
\begin{ficha}{}2 evaporadoras \md\ 1 zona \md\ Modo de etapas = Simple \md\ CON deshielo (\code{DefrostController_1})\end{ficha}

\subsection{Estructura de bloques}
\begin{arbol}
/Config/CuartosFrios/\\
`-- ColdRoom_3              (Cuarto frio)  -- Modo de etapas = Simple\\
\phantom{`-- }|-- EvaporatorUnit_1    (Evaporadora)  -- Tiene deshielo = Si\\
\phantom{`-- }|-- EvaporatorUnit_2    (Evaporadora)  -- Tiene deshielo = Si\\
\phantom{`-- }`-- DefrostController_1 (Control de deshielo)
\end{arbol}

\subsection{Configuración}
\begin{center}\small
\rowcolors{2}{white}{zebra}
\begin{tabular}{L{3.6cm} L{4.6cm} L{5.4cm}}
\rowcolor{cab}\cabfila \textbf{Bloque} & \cabfila \textbf{Slot} & \cabfila \textbf{Valor} \\
\code{ColdRoom_3} & Modo de etapas & \textbf{Simple} \\
\code{ColdRoom_3} & Consigna & consigna del cuarto \\
\code{ColdRoom_3} & Diferencial arriba / abajo & según proyecto \\
\code{EvaporatorUnit_1..2} & \textbf{Tiene deshielo} & \textbf{Sí} \\
\code{EvaporatorUnit_1..2} & Retardo de arranque & según proyecto \\
\end{tabular}
\end{center}

\noindent\textbf{Parámetros del deshielo} (\code{DefrostController_1}):
\begin{center}\small
\rowcolors{2}{white}{zebra}
\begin{tabular}{L{5.2cm} L{8.4cm}}
\rowcolor{cab}\cabfila \textbf{Slot} & \cabfila \textbf{Valor} \\
Modo & \code{Intervalo} u \code{Horario} (según proyecto) \\
Intervalo & periodo entre deshielos (si Modo = Intervalo), ej. cada 6 h \\
Entrada de horario & (si Modo = Horario) enlazar un \code{BooleanSchedule} — ver 5.5 \\
Duracion & duración máxima del deshielo \\
Retardo de escalonamiento & \textbf{4 min} (evita que ambas entren en deshielo a la vez) \\
Terminar por temp resistencia & opcional: \code{Sí} para terminar también por temperatura \\
Umbral temp resistencia & temperatura de fin de deshielo (si lo anterior = Sí) \\
\end{tabular}
\end{center}

\subsection{Entradas y salidas}
\begin{minipage}[t]{0.52\linewidth}\small
\centering\textbf{ENTRADAS}\\[4pt]
\rowcolors{2}{white}{zebra}
\begin{tabular}{L{3.0cm} L{4.2cm}}
\rowcolor{cab}\cabfila \textbf{Punto} & \cabfila \textbf{Slot destino} \\
\code{Zona1_C3} & \code{ColdRoom_3} \md \slot{Zona 1} \\
\code{Coil_C3E1} & \code{EvapUnit_1} \md \slot{Temp evap.} \\
\code{Coil_C3E2} & \code{EvapUnit_2} \md \slot{Temp evap.} \\
\code{TempResist_C3E1} & \code{EvapUnit_1} \md \slot{Temp resist.}$^\ast$ \\
\code{TempResist_C3E2} & \code{EvapUnit_2} \md \slot{Temp resist.}$^\ast$ \\
\end{tabular}
\end{minipage}\hfill
\begin{minipage}[t]{0.46\linewidth}\small
\centering\textbf{SALIDAS} (\code{in8})\\[4pt]
\rowcolors{2}{white}{zebra}
\begin{tabular}{L{4.0cm} L{2.3cm}}
\rowcolor{cab}\cabfila \textbf{Slot origen} & \cabfila \textbf{Writable} \\
\code{EvapUnit_1} \md \slot{Salida valvula} & \code{Valv_C3E1} \\
\code{EvapUnit_1} \md \slot{Salida evap.} & \code{Evap_C3E1} \\
\code{EvapUnit_1} \md \slot{Salida resist.} & \code{Resist_C3E1} \\
\code{EvapUnit_2} \md \slot{Salida valvula} & \code{Valv_C3E2} \\
\code{EvapUnit_2} \md \slot{Salida evap.} & \code{Evap_C3E2} \\
\code{EvapUnit_2} \md \slot{Salida resist.} & \code{Resist_C3E2} \\
\end{tabular}
\end{minipage}

\vspace{4pt}
\begin{nota}$^\ast$ Los sensores de temp de resistencia solo hacen falta si
\slot{Terminar por temp resistencia = Sí}. Si el deshielo termina solo por \slot{Duracion}, no los enlaces.\end{nota}

\subsection{Si el Modo del deshielo es Horario}
Crea un \code{BooleanSchedule} (WebScheduler nativo) con los horarios de deshielo del cuarto y
\textbf{enlázalo} al slot \slot{Entrada de horario} del \code{DefrostController_1}. El operador edita
esos horarios desde la vista de agenda; no se tocan los bloques.

\subsection{Cómo se comporta}
\begin{comporta}
En operación normal, ambas evaporadoras siguen a \textbf{Zona 1} (modo Simple).\\[4pt]
El \code{DefrostController_1} \textbf{coordina el deshielo}: al tocar deshelar, la secuencia por unidad
es \textbf{cerrar válvula \ra\ parar evaporadora \ra\ energizar resistencia}. Termina por \slot{Duracion}
o, si está activo, por \slot{Temp resistencia} $\geq$ \slot{Umbral}.\\[4pt]
\textbf{Interbloqueo:} nunca hay dos unidades en deshielo a la vez. Si una está deshelando y la otra ya
toca, la segunda \textbf{espera}; cuando la primera vuelve a operar, corre el \slot{Retardo de
escalonamiento} (4 min) y recién entonces la segunda inicia. La unidad que no deshiela sigue enfriando.
\end{comporta}

% ================= 6. CUARTO 4 =================
\section{Cuarto 4 — 1 evaporadora, 1 zona, sin deshielo}
\begin{ficha}{}1 evaporadora \md\ 1 zona \md\ Modo de etapas = Simple \md\ Retardo de arranque = 5 s\end{ficha}

\subsection{Estructura de bloques}
\begin{arbol}
/Config/CuartosFrios/\\
`-- ColdRoom_4           (Cuarto frio)  -- Modo de etapas = Simple\\
\phantom{`-- }`-- EvaporatorUnit_1  (Evaporadora)  -- Retardo de arranque = 5 s
\end{arbol}

\subsection{Configuración}
\begin{center}\small
\rowcolors{2}{white}{zebra}
\begin{tabular}{L{3.6cm} L{4.6cm} L{5.4cm}}
\rowcolor{cab}\cabfila \textbf{Bloque} & \cabfila \textbf{Slot} & \cabfila \textbf{Valor} \\
\code{ColdRoom_4} & Modo de etapas & \textbf{Simple} \\
\code{ColdRoom_4} & Consigna & consigna del cuarto \\
\code{ColdRoom_4} & Diferencial arriba / abajo & según proyecto \\
\code{EvaporatorUnit_1} & \textbf{Retardo de arranque} & \textbf{5 s} \\
\code{EvaporatorUnit_1} & Tiene deshielo & \textbf{No} \\
\end{tabular}
\end{center}
\begin{nota}\slot{Salida resistencia}, \slot{Temp resistencia} y \slot{Zona 2} no se usan.\end{nota}

\subsection{Entradas y salidas}
\begin{minipage}[t]{0.49\linewidth}\small
\centering\textbf{ENTRADAS}\\[4pt]
\rowcolors{2}{white}{zebra}
\begin{tabular}{L{2.5cm} L{3.9cm}}
\rowcolor{cab}\cabfila \textbf{Punto} & \cabfila \textbf{Slot destino} \\
\code{Zona1_C4} & \code{ColdRoom_4} \md \slot{Zona 1} \\
\code{Coil_C4E1} & \code{EvapUnit_1} \md \slot{Temp evap.} \\
\end{tabular}
\end{minipage}\hfill
\begin{minipage}[t]{0.49\linewidth}\small
\centering\textbf{SALIDAS} (\code{in8})\\[4pt]
\rowcolors{2}{white}{zebra}
\begin{tabular}{L{3.9cm} L{2.4cm}}
\rowcolor{cab}\cabfila \textbf{Slot origen} & \cabfila \textbf{Writable} \\
\code{EvapUnit_1} \md \slot{Salida valvula} & \code{Valv_C4E1} \\
\code{EvapUnit_1} \md \slot{Salida evap.} & \code{Evap_C4E1} \\
\end{tabular}
\end{minipage}

\subsection{Cómo se comporta}
\begin{comporta}
Modo simple con retardo mayor: cuando \textbf{Zona 1} pide frío, \textbf{abre la válvula} y
\textbf{5 s después arranca la evaporadora}. Al parar: primero la evaporadora, luego la válvula.
\end{comporta}

% ================= 7. CHECKLIST =================
\section{Lista de verificación por cuarto}
\begin{itemize}[label=$\square$,leftmargin=1.6em,itemsep=3pt]
  \item \code{ColdRoom_N} creado en \code{/Config} y renombrado.
  \item Evaporadoras arrastradas \textbf{dentro} del \code{ColdRoom_N} y numeradas \code{_1.._N}.
  \item (Cuarto 3) \code{DefrostController_1} dentro del \code{ColdRoom_3}.
  \item \slot{Modo de etapas} fijado (Cuarto 1 = Por etapas; 2/3/4 = Simple).
  \item \slot{Consigna}, \slot{Diferencial arriba}, \slot{Diferencial abajo} cargados.
  \item \slot{Retardo de arranque} por evaporadora (Cuarto 1 = 2 s; Cuarto 4 = 5 s).
  \item \slot{Tiene deshielo} = Sí solo en las evaporadoras del Cuarto 3.
  \item (Cuarto 3) parámetros de deshielo; \slot{Retardo de escalonamiento} = 4 min.
  \item Todas las ENTRADAS enlazadas (sensor \ra\ slot).
  \item Todas las SALIDAS enlazadas al writable en nivel \code{in8}.
  \item (Cuarto 3, si aplica) \slot{Entrada de horario} $\leftarrow$ \code{BooleanSchedule}.
  \item Límites de alarma cargados donde corresponde.
  \item Prueba funcional: forzar demanda de zona y verificar el orden válvula\,\ra\,evaporadora (y, en el
        Cuarto 3, la secuencia e interbloqueo de deshielo).
\end{itemize}

% ================= 8. RESUMEN I/O =================
\section{Resumen de I/O por cuarto}
\begin{center}\small
\rowcolors{2}{white}{zebra}
\begin{tabular}{L{1.6cm} L{4.4cm} L{5.0cm} L{1.8cm}}
\rowcolor{cab}\cabfila \textbf{Cuarto} & \cabfila \textbf{Entradas numéricas} & \cabfila \textbf{Salidas booleanas} & \cabfila \textbf{Deshielo} \\
1 & 2 zona + 3 serpentín = 5 & 3 válvula + 3 evaporadora = 6 & No \\
2 & 1 zona + 1 serpentín = 2 & 1 válvula + 1 evaporadora = 2 & No \\
3 & 1 zona + 2 serpentín (+2 temp resist.$^\ast$) & 2 válvula + 2 evaporadora + 2 resistencia = 6 & Sí \\
4 & 1 zona + 1 serpentín = 2 & 1 válvula + 1 evaporadora = 2 & No \\
\end{tabular}
\end{center}
\begin{nota}$^\ast$ Las 2 entradas de temp de resistencia del Cuarto 3 solo si se usa \slot{Terminar por temp resistencia}.\end{nota}

\end{document}
